\documentclass[10pt,a4paper]{article}
\usepackage[margin=22mm]{geometry}
\usepackage[T1]{fontenc}
\usepackage{lmodern,microtype,booktabs,tabularx}
\usepackage[hidelinks]{hyperref}
\newcommand{\planversion}{\input{VERSION}}
\title{HLA pangenome construction and benchmarking\\\large Technical design notes, v\planversion}
\author{Robert Hoehndorf}
\date{15 September 2026}
\begin{document}
\maketitle
\vspace{-1.5em}

\noindent\textbf{Scope.} These notes retain the detailed study design for implementation and later extensions. The three-day sprint, priorities, and decision points are specified in \texttt{plan.pdf}. The full experiment matrix and an inference prototype are sprint targets where data and cluster throughput permit; further validation follows the available evidence.

\section{Objective and design review}
We will construct a pangenome graph from more than 700 phased major
histocompatibility complex (MHC) haplotypes, including the planned Korea and CRC
additions. The study emphasises Asian diversity within a mixed-ancestry panel.
We will test whether the graph improves HLA typing from 1000 Genomes Project
(1000G) short reads relative to GRCh38 (hg38). The principal endpoint is HLA
genotype concordance with independent experimental typing. Secondary endpoints
are haplotype reconstruction, ancestry-specific performance, and computational
cost. Long-read comparisons and sequence-based classification extend this core
benchmark.

The local analysis snapshot contains 610 distinct haplotype identifiers:
106 APR, 464 HPRC release 2, 20 Japanese and 18 Saudi JaSaPaGe haplotypes,
and two references. Adding Korea and CRC is expected to bring the panel above
700 haplotypes. HPRC includes Asian and non-Asian donors; determine ancestry
per donor and quantify the final composition in an accession-level inventory.
Existing bundle plots and Immuannot
annotations provide starting materials. The current CWL mapping workflow already
performs k-mer-based haplotype sampling and produces reference-projected BAM;
this study also requires graph-native alignments. Existing annotations describe
the assemblies and will serve as graph metadata. Experimental typing supplies
the independent benchmark. These distinctions, together with donor exclusion
and matched read inputs, determine how an improvement can be attributed.

\section{Graph, samples, and ground truth}
\textbf{Graph construction.}
Define one MHC interval and homologous flanking anchors in a versioned BED file,
with explicit GRCh38 coordinates and coordinate convention. Extract phased
assembly sequence across these anchors; retain sample, donor, haplotype, cohort,
and accession provenance. Audit completeness, gaps, orientation, gene copy
number, contamination, duplicate donors across projects, and duplicated assembled
haplotypes. Distinguish supported biological homozygosity from assembly collapse.
Build the alignment graph with Minigraph--Cactus, preserving phased donor paths
and a GRCh38 path. Use existing PGGB/ODGI and pgr-tk analyses for structural
inspection. Validate sequence recovery and the topology required for Giraffe
indexing and haplotype sampling before scaling. Catalogue classical HLA genes,
secondary DRB genes, paralogues, and gene presence/absence.

\textbf{Evaluation cohort.}
Intersect 1000G sequence accessions with experimental HLA types and sample
metadata. Prefer high-coverage whole-genome data for the primary experiment;
keep legacy low-coverage data as a separate stratum. Reconcile sample aliases
and relatedness across all source projects. Exclude evaluation donors, their
alternative assemblies, and close relatives from graph construction and subset
selection. Keep development and final evaluation families separate; freeze
parameters on the development set. Where necessary, rebuild donor-excluded
graphs per evaluation fold, with identical folds for every method.

\textbf{Truth audit.}
Inventory the experimental HLA data reported as partially downloaded on the
DDBJ cluster, including paths, download completeness, source publications,
assays, loci, resolution, ambiguity, and sample overlap. This plan has reviewed
local files; cluster availability remains to be verified. The Gourraud et al.
Sanger dataset provides an experimental source for HLA-A, -B, -C, -DRB1, and
-DQB1~\cite{truth}. IGSR also distributes later HLA datasets; establish the
experimental or computational provenance of each before assigning its role~\cite{igsr}.
Freeze an IPD-IMGT/HLA release and allele-name conversion tables. Preserve
ambiguous allele sets and expression suffixes. Score each locus at the resolution
supported by its assay; use additional independently typed loci when available.

\section{Read recruitment and alignment}
Decode each CRAM with its exact reference, restoring original read orientation,
qualities, and read groups into paired FASTQ plus explicit singleton files.
Verify that unmapped reads are present. Emit each sequenced read once, avoiding
copies from secondary and supplementary alignment records. Define two inputs:
\begin{enumerate}
\item \textbf{All reads:} all recoverable reads in the complete CRAM.
\item \textbf{Prefiltered reads:} read templates with any hg38 alignment overlapping
      the declared MHC interval and flanks or relevant HLA alternate/decoy
      sequences, plus every unmapped read. Retrieve both mates for every selected
      template, including mates mapped elsewhere. Deduplicate by read group,
      query name, and mate; retain singletons explicitly.
\end{enumerate}
Use the same extraction policy across conditions and audit read counts and mate
retention. Filtering can miss HLA reads assigned elsewhere by the original
aligner; quantify this loss against the all-read condition and simulated reads
with known origins. Avoid a mapping-quality threshold during initial recruitment.

Align both inputs afresh to hg38 with BWA-MEM2 and to the HLA graph with Giraffe.
Keep the existing 1000G alignments as an additional provenance baseline.
For comparable locus specificity, embed the MHC graph in a fixed GRCh38 genomic
background, retaining appropriate paralogous and decoy sequence and replacing
the corresponding primary MHC segment. Use that same background in every graph
condition. Treat an HLA-only graph as a labelled sensitivity analysis because
whole-genome reads can map spuriously into a regional reference. Archive
Giraffe GAM alignments and their graph/index identities; derive projected BAM
for compatible downstream tools and quantify projection loss. Include Giraffe
on a GRCh38-only graph as a mapper control. Match read preprocessing and
computational resources across arms.

\section{Haplotype inference and HLA typing}
\textbf{Sequence inference.}
First infer two supported sequences per diploid HLA locus. For graph alignments,
fit candidate path pairs to read and paired-fragment likelihoods, allowing
multi-mapping and sequencing error; retain posterior ambiguity and low-coverage
masks. Use local variant calling, phasing, and assembly to recover supported
sequence absent from the donor panel. For hg38 alignments, use local variant
calling and read-backed phasing with the same coverage and confidence criteria.
Apply the same sequence-to-allele annotation procedure to both outputs. Report
per-gene haplotypes, phase-block spans, and supported inter-gene links separately;
short-read support will determine how much of the extended MHC can be phased.
Validate this inference stage on held-out simulated diplotypes before real-data
benchmarking. Model variable-copy DRB loci separately from the primary diploid
loci.

\textbf{Multiple typers.}
Benchmark HLA*LA, T1K, and SpecHLA on short reads~\cite{hlala,t1k,spechla}.
For reconstructed sequences, compare Immuannot with SpecHLA's assembly-typing
mode using compatible database releases. For genuine long-read datasets, use
SpecImmune and HLAminer, with SpecHLA as an optional additional method~\cite{specimmune,hlaminer}.
HLA*LA's long-read mode is an exploratory G-group-resolution comparator~\cite{hlala}.

Separate the common sequence-inference experiment from each typer's native
end-to-end pipeline. Many typers recruit and realign reads internally; provide
their supported FASTQ or reference-coordinate BAM inputs and explicitly identify
which upstream alignment information they retain. A FASTQ-based caller using
identical reads and its own reference runs once per read set. Compare
alignment-dependent recruitment by exporting reads and mates from each upstream
alignment. Label this as a recruitment comparison. Register a compatibility table
of input type, loci, output resolution, and internal realignment before running
the benchmark; unsupported combinations receive an explicit unavailable entry.

\textbf{Short versus long reads.}
Our hypothesis is that accurate long reads improve full-gene sequence recovery,
paralogue assignment, and phase continuity by spanning informative variants.
The gain in common two-field typing may be smaller. Read accuracy, allele balance,
coverage, and caller design determine the outcome. Compare real PacBio HiFi and,
where available, ONT data from the same independently typed individuals; analyse
technologies separately at native and matched effective HLA coverage. Report
sequencing cost alongside accuracy. Assembly-derived reads belong to simulation
experiments. Haplotype FASTA inferred from short reads remains a short-read
result. If matched long-read samples are scarce, define a separate external
validation cohort and report its composition.

\section{Comparison matrix and evaluation}
Table~\ref{tab:matrix} defines the reference axis. Cross every row with all-read
and prefiltered input. Apply compatible typers and the common sequence-inference
procedure as described above.
\begin{table}[ht]
\centering\small
\begin{tabularx}{\textwidth}{@{}lX@{}}
\toprule
Reference family & Levels and selection rule \\
\midrule
Linear controls & hg38/BWA-MEM2; GRCh38-only graph/Giraffe. \\
Full HLA graph & All eligible donor haplotypes after evaluation exclusions. \\
Whole-donor subsets & 4, 8, 16, 32 individuals, retaining both available phased
haplotypes (nominally 8, 16, 32, 64 paths). Compare read-k-mer-ranked donors with
ancestry-stratified random subsets; use five fixed random seeds. \\
Blockwise personalisation & 4, 8, 16, 32 haplotypes per block selected from read
k-mers; optional diploid selection from 32 candidates. \\
\bottomrule
\end{tabularx}
\caption{Planned reference conditions. Reference/background paths are additional
to the stated donor or haplotype counts.}\label{tab:matrix}
\end{table}
The cited personalisation method samples \emph{haplotypes per approximately
10-kb block}, and assumes adequate k-mer coverage (at least 20-fold in its
design)~\cite{personal}. Whole-individual subsetting is a separate experiment.
Select donors using sequence evidence with a fixed scoring rule and tie-breaker.
For the main matrix, derive selection k-mers from the same input read set used
for mapping. A secondary experiment holds all-read k-mer selection fixed while
varying mapping input to separate selection loss from mapping loss. Freeze
blocks, seeds, donor lists, indexes, and coverage rules. Final-test truth labels
remain hidden during selection and tuning.

\textbf{Primary analysis.}
Pre-specify full graph versus hg38 on all reads as the main contrast. Measure
exact unordered diploid genotype concordance at two-field resolution where
experimentally resolved, with a separate G-group analysis where that is the
available resolution. Include no-calls in the eligible denominator. Report
call rate and accuracy among calls, allele-copy accuracy, and ambiguity-compatible
concordance with predicted-set size. Evaluate each locus and a macro-average
over a fixed shared locus set. Report all supported loci separately for broader
typers. Resolve nomenclature differences before scoring; preserve uncertain
experimental labels as sets.

\textbf{Secondary analysis.}
Stratify by ancestry, graph representation, allele frequency, coverage, and
structural complexity. Measure filtering retention, locus misassignment in
simulation, allelic balance, sequence error over callable bases, phase-block
completeness, and switch error where phased truth exists. Experimental HLA
labels assess typing; independent phased sequence is required to assess extended
haplotype reconstruction. Report wall time, CPU time, peak memory, and storage,
including extraction, selection, indexing, and typing. Use paired comparisons
and confidence intervals from resampling individuals or families; pre-specify
multiplicity control for secondary tests. Establish the required test-set size
from development-set discordance and a declared minimum useful improvement.

\section{Exploratory sequence and functional classification}
Begin with label-blind, per-locus clustering of principal bundles: recurrent
sequence segments shared along haplotype paths. Represent each haplotype by
bundle presence, copy number, order, and intervening sequence. Compare a
length-weighted bundle distance with nucleotide and peptide distances. Analyse
whole genes, peptide-binding exons, and extended class II separately to distinguish
coding similarity from linkage and structural variation. Select clustering
parameters by stability under resampling and graph-parameter perturbation;
then test recovery of allele groups, two-field types, and DRB structures on
held-out donors using adjusted Rand index and adjusted mutual information.
The existing pgr-tk bundle analyses provide a feasibility starting point.

Prioritise functional interpretation through translated binding-site residues
and peptide-binding profiles on a fixed peptide panel, followed by independent
experimental binding evidence where available. Test whether sequence clusters
add information beyond named HLA alleles. Population frequencies provide
descriptive context with sample-size uncertainty and ancestry stratification.
Phenotype associations require individual-level outcomes, ancestry and kinship
adjustment, comparison against conventional HLA predictors, multiple-testing
control, and external replication. Population disease incidence can motivate
hypotheses; individual-level cohorts provide the association test.

\section{Execution and deliverables}
First reconcile the assembly and truth manifests, then complete a small,
ancestry-diverse development pilot spanning all-read and prefiltered inputs,
hg38, and the full graph. Verify graph integrity, donor exclusion, recruitment,
typing interfaces, and resource requirements before expanding to the subset
matrix. Freeze the evaluation cohort and analysis specification before final
benchmarking. Deliver versioned manifests and configuration, reproducible CWL
workflows, graphs and indexes, per-sample calls with uncertainty, benchmark
summaries, and a separate exploratory classification analysis. Pin tools,
containers, database releases, seeds, and source checksums in machine-readable
provenance. This document specifies proposed work and expected tests.

\begin{thebibliography}{9}\small
\bibitem{personal} Sir\'en J et al. Personalized pangenome references.
\emph{Nature Methods} 21, 2017--2023 (2024).
\href{https://doi.org/10.1038/s41592-024-02407-2}{doi:10.1038/s41592-024-02407-2}.
\bibitem{truth} Gourraud P-A et al. HLA diversity in the 1000 Genomes dataset.
\emph{PLOS ONE} 9, e97282 (2014).
\href{https://doi.org/10.1371/journal.pone.0097282}{doi:10.1371/journal.pone.0097282}.
\bibitem{igsr} IGSR. \href{https://www.internationalgenome.org/faq/was-hla-diversity-studied-in-igsr/}{Was HLA diversity studied in IGSR?}
\bibitem{hlala} Dilthey Lab. \href{https://github.com/DiltheyLab/HLA-LA}{HLA*LA documentation};
\href{https://pmc.ncbi.nlm.nih.gov/articles/PMC6821427/}{HLA typing from linearly projected graph alignments} (2019).
\bibitem{t1k} Song L et al. Efficient and accurate KIR and HLA genotyping with massively parallel sequencing data (2023).
\href{https://doi.org/10.1101/gr.277585.122}{doi:10.1101/gr.277585.122}; \href{https://github.com/mourisl/T1K}{T1K documentation}.
\bibitem{spechla} DeepOmics Lab. \href{https://github.com/deepomicslab/SpecHLA}{SpecHLA documentation}.
\bibitem{specimmune} DeepOmics Lab. \href{https://github.com/deepomicslab/SpecImmune}{SpecImmune documentation}.
\bibitem{hlaminer} Birol Lab. \href{https://github.com/BirolLab/HLAminer}{HLAminer documentation}.
\end{thebibliography}
\noindent\footnotesize Web documentation consulted 15 September 2026. Software versions will be frozen in the execution manifest.
\end{document}
